\chapter{Notation and Game Recording}

RPSC의 기보법은 체스의 algebraic notation과 PGN의 기록 방식을 가능한 한 유지하면서, RPSC에만 필요한 퀴즈 결과, 주사위 이동 경로, 아이템을 추가한다. 기보의 목적은 경기의 흐름을 간결하게 읽는 동시에, 필요한 경우 각 말의 위치와 방향을 처음부터 다시 복원할 수 있게 하는 것이다. 실제 경기의 기록과 분석용 표시는 구분하되, 둘 모두 같은 Move Notation을 사용한다.

\section{Board and Pieces}

보드의 좌표는 체스와 같은 algebraic coordinates를 사용한다. White 쪽에서 보았을 때 왼쪽 아래 칸을 \texttt{a1}로 두고, 오른쪽으로 \texttt{a}부터 \texttt{h}, 위쪽으로 \texttt{1}부터 \texttt{8}을 붙인다. 체스의 통상적인 색칠을 가정하면 \texttt{a1}은 어두운 칸이다. 기보에서는 보드를 White가 아래쪽, Black이 위쪽에 오도록 돌려 기록한다.

최초 선공 \emph{보드 역할}을 \textbf{White}, 최초 후공 보드 역할을 \textbf{Black}이라 부른다. White와 Black은 학교의 고정된 이름이 아니며, \POSTECH{}과 \KAIST{} 역시 White/Black의 별칭이 아니다. 경기마다 \texttt{WhiteTeam}과 \texttt{BlackTeam}에 두 학교의 배정을 기록한다. 반면 \texttt{White}와 \texttt{Black} header에는 그 보드 역할을 실제로 조종한 선수, 선수단 또는 Engine의 이름을 기록한다. 각 보드 역할의 네 말에는 경기 내내 변하지 않는 고유 번호를 붙인다.

\begin{center}\renewcommand{\arraystretch}{1.15}\begin{tabular}{@{}cccc@{}}\toprule
말 & 최초 위치 & 최초 윗면 & 번호 기준 \\\midrule
W1 & a1 & S & White의 왼쪽부터 첫 번째 \\
W2 & c1 & R & White의 왼쪽부터 두 번째 \\
W3 & e1 & P & White의 왼쪽부터 세 번째 \\
W4 & g1 & S & White의 왼쪽부터 네 번째 \\
B1 & h8 & S & Black의 왼쪽부터 첫 번째 \\
B2 & f8 & R & Black의 왼쪽부터 두 번째 \\
B3 & d8 & P & Black의 왼쪽부터 세 번째 \\
B4 & b8 & S & Black의 왼쪽부터 네 번째 \\\bottomrule
\end{tabular}\end{center}

\begin{figure}[htbp]
    \centering
    \begin{tikzpicture}[x=0.72cm,y=0.72cm]
        % Board: line-only, with standard algebraic coordinates.
        \foreach \i in {0,...,8} {
            \draw[rpsc board line] (\i,0)--(\i,8);
            \draw[rpsc board line] (0,\i)--(8,\i);
        }
        \draw[rpsc board border] (0,0) rectangle (8,8);

        \foreach \x/\file in {0.5/a,1.5/b,2.5/c,3.5/d,4.5/e,5.5/f,6.5/g,7.5/h} {
            \node[rpsc coordinate] at (\x,-0.33) {\file};
        }
        \foreach \y/\rank in {0.5/1,1.5/2,2.5/3,3.5/4,4.5/5,5.5/6,6.5/7,7.5/8} {
            \node[rpsc coordinate] at (-0.32,\y) {\rank};
        }

        % White: upright from White's seated direction. Numbered from White's left.
        \rpscglyph{0.5}{0.58}{S}{0}{\RPSCWhiteColor}
        \rpscglyph{2.5}{0.58}{R}{0}{\RPSCWhiteColor}
        \rpscglyph{4.5}{0.58}{P}{0}{\RPSCWhiteColor}
        \rpscglyph{6.5}{0.58}{S}{0}{\RPSCWhiteColor}
        \node[rpsc piece id] at (0.5,0.18) {W1};
        \node[rpsc piece id] at (2.5,0.18) {W2};
        \node[rpsc piece id] at (4.5,0.18) {W3};
        \node[rpsc piece id] at (6.5,0.18) {W4};

        % Black: upright from Black's seated direction, hence rotated 180 degrees in White view.
        \rpscglyph{7.5}{7.42}{S}{180}{\RPSCBlackColor}
        \rpscglyph{5.5}{7.42}{R}{180}{\RPSCBlackColor}
        \rpscglyph{3.5}{7.42}{P}{180}{\RPSCBlackColor}
        \rpscglyph{1.5}{7.42}{S}{180}{\RPSCBlackColor}
        \node[rpsc piece id] at (7.5,7.82) {B1};
        \node[rpsc piece id] at (5.5,7.82) {B2};
        \node[rpsc piece id] at (3.5,7.82) {B3};
        \node[rpsc piece id] at (1.5,7.82) {B4};
    \end{tikzpicture}
    \caption{기보에서 사용하는 표준 방향과 초기 배치. 각 팀의 말은 해당 팀이 자기 자리에서 보았을 때 정방향으로 표시한다. 이 예시에서는 \POSTECH{}이 White, \KAIST{}가 Black이다.}
    \label{fig:initial-position}
\end{figure}


보드 그림에서 말은 $\bar{\mathbf{S}}$, $\bar{\mathbf{R}}$, $\bar{\mathbf{P}}$로 표시한다. bar는 문자 바로 위에 놓이며, 문자와 bar 전체를 하나의 기호로 취급하여 함께 회전시킨다. 특히 S는 문자만으로는 $180^\circ$ 회전을 구별하기 어려우므로 bar를 통해 방향을 명확하게 나타낸다.

\section{Move Notation}

말의 일반 이동은 ``말 번호: 이동 경로''의 순서로 적는다. 콜론 뒤에는 한 칸을 띄운다. 이동 경로에는 \textbf{그 수를 시작하기 직전의 출발칸을 먼저 적고}, 이후 각 Roll로 도달한 모든 칸을 시간 순서대로 하이픈으로 연결한다.
\begin{quote}\ttfamily W2: c1-c2-d2-d3-e3\end{quote}
출발칸과 도착칸만이 아니라 모든 경유칸을 기록하므로, 최초 주사위 방향을 알고 있으면 기보만으로 이동 후의 윗면과 방향까지 복원할 수 있다.

아이템을 사용한 이동은 말 번호 뒤의 대괄호 안에 아이템 코드를 적는다. 아이템은 첫 번째 Roll보다 먼저 적용한다.
\begin{center}\renewcommand{\arraystretch}{1.15}\begin{tabular}{@{}cL{0.66\textwidth}@{}}\toprule
코드 & 의미 \\\midrule
\texttt{Pu} & Push \\
\texttt{RoN} & North로 한 칸 Roll했을 때와 같은 orientation 변화를 현재 칸에서 적용 \\
\texttt{RoS} & South로 한 칸 Roll했을 때와 같은 orientation 변화를 현재 칸에서 적용 \\
\texttt{RoE} & East로 한 칸 Roll했을 때와 같은 orientation 변화를 현재 칸에서 적용 \\
\texttt{RoW} & West로 한 칸 Roll했을 때와 같은 orientation 변화를 현재 칸에서 적용 \\
\texttt{RoL} & 위에서 보아 반시계방향으로 $90^\circ$ 제자리 회전 \\
\texttt{RoR} & 위에서 보아 시계방향으로 $90^\circ$ 제자리 회전 \\
\texttt{StL} & Step Long \\
\texttt{StS} & Step Short \\\bottomrule
\end{tabular}\end{center}

Rotation의 N, S, E, W는 White/Black의 시점과 무관한 \textbf{보드 절대방향}이다. North는 rank 증가, South는 rank 감소, East는 file 증가, West는 file 감소를 뜻한다. \texttt{RoN}, \texttt{RoS}, \texttt{RoE}, \texttt{RoW}는 해당 방향으로 한 칸 Roll했을 때의 정육면체 회전만 제자리에서 적용한다. \texttt{RoL}, \texttt{RoR}은 보드에 수직인 축을 중심으로 회전한다.

초기 W2를 예로 들면 Rotation 후 윗면에 따라 Roll 수가 달라질 수 있으므로 다음과 같이 기록된다.
\begin{quote}\ttfamily
W2[RoN]: c1-c2-d2-d3\\
W2[RoE]: c1-c2-d2-d3-e3-e4\\
W2[RoR]: c1-c2-d2-d3-e3
\end{quote}

Push로 밀린 한 칸은 Roll과 구별하기 위해 \texttt{>}로 적는다. 첫 Roll은 Push 이전 칸으로 되돌아갈 수 있다. 아래 예시는 첫 Roll 직전 윗면이 P인 말의 다섯 Roll을 나타낸다. 이후 연속된 Roll끼리의 즉시 역행은 여전히 금지된다.
\begin{quote}\ttfamily W3[Pu]: e3>f3-e3-e4-f4-f5-g5\end{quote}

아이템을 새로 획득한 경우에는 \textbf{학교명이 아니라 보드 역할}의 팀 문자 뒤에 \texttt{+}와 기본 아이템 코드를 붙인다. Rotation과 Step은 획득 시에는 방향이나 길이를 아직 선택하지 않으므로 각각 \texttt{Ro}, \texttt{St}만 기록한다.
\begin{quote}\ttfamily W+Pu\qquad B+Ro\qquad W+St\end{quote}

제한 시간 초과는 그 결과로 실행되는 행동의 \textbf{앞}에 \texttt{Timeout}을 적는다. 기보는 실제 사건이 일어난 순서를 따른다. 이동 시간 초과에서도 실제로 나온 full path를 모두 기록하고, 아이템 선택 시간 초과에서는 실제로 무작위 획득된 아이템을 기록한다.
\begin{quote}\ttfamily
Timeout W2: c1-c2-d2-e2-e3\\
Timeout W+Ro
\end{quote}

전투 결과 말이 제거되면 이동 기보 뒤에 \texttt{x}와 제거된 말의 번호를 적는다. \texttt{x}는 ``뒤의 말이 제거되었다''는 뜻이다. 한 팀의 네 말이 모두 제거되어 초기화가 발생하면 마지막에 \texttt{Reset}을 붙인다.
\begin{quote}\ttfamily B1: h5-g5-g4-f4 xW4 Reset\end{quote}

따라서 한 행동 안의 정보는 필요한 경우
\[\texttt{Timeout}\;\longrightarrow\;\text{Move or Item Gain}\;\longrightarrow\;\text{Capture}\;\longrightarrow\;\texttt{Reset}\]
의 실제 발생 순서로 적는다. 이동 경로는 경기 복원의 기준이 되는 \textbf{canonical record}이다.

\section{Game Record}

기보의 기본 단위는 보드의 한 수가 아니라 \textbf{과학퀴즈 한 문제}이다. 각 줄의 앞에 문제 번호를 적고, 바로 뒤에
\[
\boxed{\texttt{Q[POSTECH, KAIST]}}
\]
순서로 두 학교의 정오 여부를 기록한다. 정답은 1, 오답은 0으로 적으며 쉼표 뒤에는 한 칸을 띄운다. 이 순서는 White/Black 배정과 무관하게 모든 경기에서 고정된다.

\begin{center}\renewcommand{\arraystretch}{1.12}\begin{tabular}{@{}cc@{}}\toprule
표기 & 의미 \\\midrule
\texttt{Q[1, 1]} & 양교 모두 정답 \\
\texttt{Q[1, 0]} & \POSTECH{}만 정답 \\
\texttt{Q[0, 1]} & \KAIST{}만 정답 \\
\texttt{Q[0, 0]} & 양교 모두 오답 \\\bottomrule
\end{tabular}\end{center}

한 학교만 정답이면 먼저 \texttt{Q}로 정답 학교를 결정하고, 그 학교의 현재 보드 역할을 두 header \texttt{WhiteTeam}과 \texttt{BlackTeam}에서 확인하여 \texttt{W+...} 또는 \texttt{B+...}를 적는다. 예를 들어 역배정 경기에서는 다음이 canonical record이다.
\begin{Verbatim}[fontsize=\small]
[WhiteTeam "KAIST"]
[BlackTeam "POSTECH"]

Q[1, 0] B+St
Q[0, 1] W+Ro
\end{Verbatim}
즉 \texttt{Q}는 학교를 보고, \texttt{W/B}는 보드를 본다.

양교의 정오 여부가 같으면 White의 이동을 먼저, Black의 이동을 다음에 적는다. 다음 예시는 WhiteTeam=POSTECH, BlackTeam=KAIST인 경우이다.
\begin{quote}\ttfamily
1. Q[1, 1] W2: c1-c2-d2-d3-e3 B2: f8-f7-e7-e6-d6\\
2. Q[1, 0] W+Ro\\
3. Q[0, 1] B+St\\
4. Q[0, 0] W1[StL]: a1-a2-b2-b3-c3 B3[RoR]: d8-d7-c7-c6-b6-b5
\end{quote}

경기 전체 기록은 PGN처럼 header와 game record로 나눈다. \texttt{White}/\texttt{Black}은 플레이어 또는 controller, \texttt{WhiteTeam}/\texttt{BlackTeam}은 학교 배정을 뜻한다. \texttt{Result}, \texttt{Score}, \texttt{Quiz}, \texttt{Captures}의 두 수는 모두 White--Black 순서이다. 이들과 달리 각 문제의 \texttt{Q[...]}만 POSTECH--KAIST 순서로 고정된다.

\begin{Verbatim}[fontsize=\small]
[Format "3"]
[Event "2026 KAIST-POSTECH Science War, Science Quiz"]
[Site "POSTECH"]
[Date "2026.09.09"]
[White "Engine"]
[Black "POSTECH Squad"]
[WhiteTeam "KAIST"]
[BlackTeam "POSTECH"]
[Result "1-0"]
[Score "25-23"]
[Quiz "11-11"]
[Captures "7-6"]
[QOrder "POSTECH, KAIST"]
\end{Verbatim}

\texttt{Result}는 White 승리 \texttt{1-0}, Black 승리 \texttt{0-1}, 아직 결과가 정해지지 않은 경기 \texttt{*}를 사용한다. 본 문제는 \texttt{1.}부터 \texttt{20.}까지 기록한다. 총점과 퀴즈 정답 수가 모두 같아 연장전이 진행되면 \texttt{OT1.}, \texttt{OT2.}, \ldots{} 형식으로 퀴즈 결과만 기록한다. Analysis Board의 Human vs Human, Human vs Engine, Engine vs Engine 모드도 모두 같은 Quiz 진행과 승패 판정 구조를 따른다.

\section{Analysis}

실제 경기의 Game Record에서는 \texttt{1.}, \texttt{2.}, \ldots{}가 과학퀴즈 문제 번호이다. 반면 퀴즈의 정오 여부와 무관하게 보드의 수순 자체를 설명할 때에는 가상의 \texttt{Q[...]}를 넣지 않고 별도의 \textbf{Analysis Line}을 사용한다. 수순 번호에는 \texttt{M}을 붙이며 이를 \textbf{M-number}라 부른다.
\begin{quote}\ttfamily
M1. W2: c1-c2-d2-d3-e3 B2: f8-f7-e7-e6-d6\\
M2. W1: a1-a2-b2-b3 B3[RoR]: d8-d7-c7-c6-b6-b5
\end{quote}
White의 수만 적을 때에는 \texttt{Mn.}, Black의 대응만 적거나 Black의 수에서 variation을 시작할 때에는 \texttt{Mn...}을 사용한다. 한 팀이 보드에서 한 번 움직인 것을 \textbf{ply}라 부르므로, 완전한 M-number 하나는 White와 Black의 두 ply로 이루어진다.

\medskip\noindent\textbf{Move Annotations.}
\begin{center}\renewcommand{\arraystretch}{1.15}\begin{tabular}{@{}cL{0.69\textwidth}@{}}\toprule
기호 & 의미 \\\midrule
\texttt{!!} & 매우 뛰어나고 찾기 어려운 수 \\
\texttt{!} & 좋은 수, 특히 중요한 정확한 수 \\
\texttt{!?} & 흥미롭지만 평가가 완전히 확정되지 않은 수 \\
\texttt{?!} & 의심스럽거나 위험 요소가 있는 수 \\
\texttt{?} & 명확한 실수 \\
\texttt{??} & 평가를 크게 악화시키는 중대한 실수 \\\bottomrule
\end{tabular}\end{center}

\medskip\noindent\textbf{Position Evaluation.}
Position Evaluation은 White의 관점에서 적는다. $=$, $+=$, $=+$, $\pm$, $\mp$, $+-$, $-+$, $\infty$는 각각 균형, White/Black의 약간의 우세, 명확한 우세, 결정적 우세, 불명확을 나타내는 전통적인 체스식 분석 표기이다.

\medskip\noindent\textbf{Incomplete Records and Analysis Sessions.}
Analysis Board의 \texttt{.rpsc} 파일은 경기 종료 시점뿐 아니라 진행 중인 확정 상태에서도 저장할 수 있다. 같은 정오 결과의 \texttt{Q[1,1]} 또는 \texttt{Q[0,0]}만 기록되어 있으면 그 문제의 White 이동 전 상태를 뜻하고, 그 뒤 White의 Move Notation만 있으면 Black 이동 전 상태를 뜻한다. 한 학교만 정답인 경우에는 \texttt{Q}가 가리키는 정답 학교의 Item Choice가 아직 남아 있을 수 있다.

RPSC 0.20.0부터 Analysis Board가 저장하는 \texttt{.rpsc}는 \textbf{Format 3}을 사용한다. Format 3의 Main Line에서 \texttt{Q}는 언제나 POSTECH--KAIST 순서이고, \texttt{White}와 \texttt{Black}은 controller, \texttt{WhiteTeam}과 \texttt{BlackTeam}은 학교 배정이다. Safety tag \texttt{QOrder}의 값은 \texttt{POSTECH, KAIST}로 고정하여 사람이 파일을 직접 열어도 이 의미를 즉시 확인할 수 있게 한다.

기존 \textbf{Format 2} 파일은 호환을 위해 계속 불러올 수 있다. Format 2의 \texttt{White}/\texttt{Black}은 학교 배정이며, 본문의 \texttt{Q[a,b]}도 legacy White--Black 순서로 해석한다. Loader는 이를 먼저 POSTECH/KAIST 내부 상태로 변환한 뒤 replay하며, 다시 저장하면 Format 3의 canonical school-order \texttt{Q}로 migration한다. Format 2의 의미를 새 규칙으로 소급해 바꾸어 읽지 않는다.

Session metadata는 복원 편의를 위한 정보이며 canonical Game Record보다 우선하지 않는다. 외부 편집기로 Main Line을 수정하여 metadata와 본문이 달라지면 Analysis Board는 본문을 처음부터 replay한다. \texttt{Score}, \texttt{Quiz}, \texttt{Captures} header는 White--Black 순서이고, 본문 replay 결과와 충돌하면 현재 상태와 수치는 본문을 기준으로 다시 계산한다.

RPSC의 수치 Evaluation은 체스의 centipawn이 아니라 \textbf{RPSC score unit}을 기준으로 한다. 공식 점수에 맞추어 이미 확정된 Quiz 정답은 $1.00$, capture 하나는 $2.00$으로 반영하며, 양수는 White, 음수는 Black의 우세를 나타낸다. 남은 본 문제에서는 Symmetric Quiz Assumption을 사용할 수 있으므로 미확정 Quiz로 점수차나 아이템을 새로 만들지 않는다.
